\section{詳細解説}

\subsection{ZMP の数値例題}

\begin{example}[静止状態の ZMP]
荷物を前方の地面から持ち上げる直前の姿勢（肩 (0, 0.75) m、上腕 $-20^\circ$、前腕は上腕から $-50^\circ$）で静止しているとき、
加速度がゼロなので ZMP は重心の水平位置に一致する。各部の水平位置は、台車（80 kg）0 m、上腕の重心（4 kg）0.235 m、前腕の重心（3 kg）0.556 m、荷物（10 kg）0.641 m なので
\[
p = \frac{4\times0.235+3\times0.556+10\times0.641}{97} = 0.093\ \text{m}.
\]
接地範囲の半分 0.4 m に対して余裕は十分ある。図~\ref{fig:sim06_tipover.png} の左上で、どの動作時間でも ZMP が 0.09 m から始まっているのはこのためである。
\end{example}

\subsection{フルード則の確認}

安全な最短動作時間を $T_{\min} = 1.616\sqrt{s\,g_{\text{地球}}/g}$ として、大きさ・重力の違う 4 条件で ZMP を大きさ $s$ で割って重ねると、
4 本の曲線は完全に一致した（図~\ref{fig:D8_froude_collapse.png}）。
付録~\ref{app:zmp} のとおり、ZMP/$s$ は無次元量 $s/(gT^2)$ だけの関数なので、これを揃えれば運動は相似になる。
\textbf{18 m 級の機体の地上での動作は、人間大の機体の動作を時間方向に $\sqrt{10}\approx3.2$ 倍に引き伸ばしたものとして、そのまま予測できる}。

\DFIG{D8_froude_collapse.png}{フルード則：動作時間を $\sqrt{s/g}$ に比例させると、大きさ・重力が違っても規格化した ZMP が 1 本の曲線に重なる。}

\DTAB{sim06_scaling}

\subsection{軌道最適化の定式化}

Part B の軌道最適化は、次の問題を SLSQP で解いた。
\begin{itemize}
\item \textbf{変数}：動作時間 $T$ と、関節ごとに 4 つの補正係数 $c_{jk}$（計 9 個）。
\item \textbf{軌道}：$q_j(\tau) = q_j^{\text{最小躍度}}(\tau) + \sum_{k=1}^4 c_{jk}\,64\,\tau^3(1-\tau)^3\sin(k\pi\tau)$。
  補正の基底関数は $\tau=0,1$ で値・速度・加速度がゼロなので、始点と終点の姿勢・静止条件を保つ。
\item \textbf{目的}：$T$ の最小化。
\item \textbf{制約}：時間方向の 2 点ごとに、転倒（ZMP が接地範囲の内側 5 cm 以内）、滑り（$|F_x|\le\mu N$）、各関節の速度 6 rad/s・加速度 60 rad/s$^2$。
\end{itemize}
初期値を 6 通り（最小躍度＋ランダムな補正）変えて解き、制約を満たす最良解を採用した。

\subsection{スティックスリップの詳細}

\DTAB{sim06_yaw}

月面で、摩擦係数と速い振りの時間を変えて 1 往復あたりのヨーを調べた（図~\ref{fig:D9_stickslip_map.png}）。
\begin{itemize}
\item 振りが速い（0.3 s）と、摩擦係数によらず約 $-50^\circ$/往復回る。慣性の反作用が摩擦の上限を大きく超え、ほぼ自由に滑るためである。
\item 振りが遅いほど、回転は小さくなり、やがて滑らなくなる。滑り始める境界は、摩擦係数が大きいほど速い振り側に移る。
\item 摩擦が小さく（$\mu\approx0.2$）振りが遅い（約 2 s）場合には、回転の向きが逆転する（$+15^\circ$ 程度）。
  ゆっくり戻すとき（3 s）にも滑るようになり、戻りの滑りの方が大きくなるためである。
\end{itemize}

\DFIG{D9_stickslip_map.png}{月面で、腕の水平振り 1 往復による台車のヨー [°]。横軸は速い振りの時間（戻りは 3 s）、縦軸は摩擦係数。}

\section*{考察課題}
\begin{exercise}
図~\ref{fig:D8_froude_collapse.png} の相似則が成り立つには、関節の速度・加速度の上限も相似に拡大されている必要がある。上限が大きさによらず一定だと、何が変わるか。
\end{exercise}
\begin{exercise}
月面のローバーにアームを載せるとき、作業中に機体が回らないようにするための設計・運用上の対策を、図~\ref{fig:D9_stickslip_map.png} から考えよ。
\end{exercise}
